\documentclass[11pt]{article}

% --- Required style ---
\usepackage{series}

% --- Math packages ---
\usepackage{amsmath,amssymb,amsthm}
\usepackage{geometry}
\usepackage{hyperref}
\usepackage{booktabs}

\geometry{margin=1in}


% --- Metadata ---
\title{Theta Closure in Modal Triplet Theory I:
Gauge Couplings from Internal Geometry}
\author{Peter Nero}
\date{January 2026}

\begin{document}
\maketitle
\begin{abstract}
We present a conservative $\Theta$--closure framework in Modal Triplet Theory
(MTT) relating Standard Model gauge couplings to internal geometric overlap
integrals.
Starting from experimentally measured couplings at the electroweak scale, we
derive explicit numerical targets for the nonabelian overlap ratios
$I_2/I_1$ and $I_3/I_1$ at a common matching scale
$\mu_\Theta=5~\mathrm{TeV}$ using one--loop renormalization group evolution.
We formulate precise admissibility and spectral gap conditions on the internal
geometry and prove the existence of a nonempty $\Theta$--region consistent with
these constraints.
To address the normalization of nonabelian harmonics, we introduce a
high--coherence (twistor) corner that provides canonical massless representatives
with controlled $O(\lambda_Q^{-1})$ errors.
The present paper establishes the quantitative framework and targets of
$\Theta$--closure; explicit geometric realization and independent cross--checks
are carried out in companion papers.
\end{abstract}

\section{Purpose and scope}

This paper presents a rigorous, explicit derivation linking Standard Model gauge couplings
to internal geometric data in Modal Triplet Theory (MTT).

The goals are deliberately narrow:
\begin{enumerate}
\item Derive gauge–coupling expressions from internal overlap integrals.
\item Extract numerical $\Theta$–constraints from experimentally measured couplings.
\item Prove existence of a nonempty admissible $\Theta$–region consistent with MTT gap conditions.
\item State all assumptions and calibration choices explicitly.
\end{enumerate}

No claim of uniqueness or full unification is made.
All results are conditional and falsifiable.

\section{Standing assumptions of Modal Triplet Theory}

We collect here the assumptions used throughout.
They are standard within the MTT corpus and are not re-derived.

\begin{assumption}[Bounded internal geometry]
The internal bundles $B_n$ possess bounded geometry and admit commuting vertical Laplacians.
\end{assumption}

\begin{assumption}[Uniform spectral gap]
There exists a strictly positive lower bound
\[
\lambda_{\ast} > 0
\]
separating coherent from noncoherent internal modes.
\end{assumption}

\begin{assumption}[Bounded coherent projection]
The joint Riesz projector
\[
\Pi = \Pi_{B_1}\Pi_{B_2}\Pi_{B_3}
\]
is bounded on $H^1$.
\end{assumption}

\begin{assumption}[Fundamental Contractivity Condition (FCC)]
The projected flow $F=\Pi\circ\Phi_\tau$ is contractive on the coherent sector.
\end{assumption}

\begin{remark}
If any of these assumptions fail, MTT predicts breakdown of physics
(probabilities, records, or observers cease to exist).
\end{remark}

\subsection{Spectral gap as a selection principle}

A uniform spectral gap is assumed throughout this work as a condition for the
existence of a stable coherent sector.
This assumption is not introduced for mathematical convenience alone, but
reflects a physical selection criterion.

If the spectral gap closes, the separation between coherent and noncoherent
modes is lost.
In this regime:
\begin{itemize}
\item the projection onto a finite coherent sector becomes ill-defined,
\item probabilities and amplitudes cannot be assigned consistently,
\item effective locality and observer--accessible records cease to exist.
\end{itemize}

Thus the spectral gap functions as a selection principle rather than a tunable
parameter: theories without such a gap do not admit stable effective physics.
The present work does not derive the existence of the gap dynamically; instead,
it explores the consequences of assuming a coherent sector exists.
Determining whether such a gap is inevitable, emergent, or contingent defines a
distinct and deeper problem beyond the scope of this paper.

\section{Internal bundle structure}

MTT employs a tri-bundle internal structure.
At each spacetime point $y\in Y^4$,
\[
B_n|_y \simeq S^1_{\mathrm{cen}} \times \Sigma_n,
\qquad n=1,2,3.
\]

Here:
\begin{itemize}
\item $S^1_{\mathrm{cen}}$ is a common central circle,
\item $\Sigma_1$ is an auxiliary circle,
\item $\Sigma_2$ is a lens-type surface,
\item $\Sigma_3$ is a nilmanifold leaf.
\end{itemize}

This structure is fixed throughout the paper.

\section{Gauge couplings from internal overlaps}

\subsection{Definition of overlap integrals}

\begin{definition}[Overlap integrals]
For each gauge factor $a=1,2,3$, define
\[
I_a := \int_{B_a} \langle \omega_a, \omega_a \rangle \, d\mu_{B_a},
\]
where $\omega_a$ is the normalized harmonic representative on $B_a$.
\end{definition}

\subsection{Reduction of the Yang--Mills action}

\begin{proposition}[Coupling–overlap relation]
Under canonical normalization,
\[
\frac{1}{g_a^2} = \frac{1}{g_{10}^2}\, I_a.
\]
\end{proposition}

\begin{proof}
Start from the ten-dimensional Yang–Mills action
\[
S_{10} = \frac{1}{2g_{10}^2}\int \mathrm{Tr}(F_{AB}F^{AB})\, d\mathrm{vol}_{10}.
\]

Expand the gauge field along normalized internal harmonic modes
\[
A_\mu(x,u) = A_\mu^{(a)}(x)\,\omega_a(u),
\qquad \|\omega_a\|_{L^2(B_a)}=1.
\]

Integrating over the internal space yields
\[
S_4 = \frac{1}{4g_a^2}\int F_{\mu\nu}^{(a)}F^{(a)\mu\nu}\, d\mathrm{vol}_4,
\]
with
\[
\frac{1}{g_a^2} = \frac{1}{g_{10}^2}\int_{B_a}\langle \omega_a,\omega_a\rangle d\mu_{B_a}.
\]
\end{proof}

\subsection{Normalization conventions}

We adopt the following conventions:

\begin{enumerate}
\item \textbf{Generator normalization:}
\[
\mathrm{Tr}(T^a T^b)=\frac{1}{2}\delta^{ab}.
\]

\item \textbf{Hypercharge normalization (GUT-normalized):}
\[
g_1=\sqrt{\frac{5}{3}}\, g',
\qquad
\alpha_1=\frac{5}{3}\alpha_Y.
\]
\end{enumerate}

With these conventions, the normalization constants satisfy
\[
N_1=N_2=N_3=1,
\]
and all group-theoretic factors are absorbed into the definition of $g_1$.

\section{Summary of this section}

At this point we have established:

\begin{itemize}
\item Gauge couplings descend from internal overlap integrals.
\item The relation $1/g_a^2 \propto I_a$ is exact under stated assumptions.
\item Ratios of couplings are ratios of geometric overlaps.
\end{itemize}

In the next part we will:
\begin{enumerate}
\item Compute $g_1,g_2,g_3$ explicitly from experimental inputs.
\item Run them to a common matching scale.
\item Extract numerical $\Theta$–constraints.
\end{enumerate}
% ===========================
\section{Experimental inputs and gauge couplings at $M_Z$}
% ===========================

In this section we compute the Standard Model gauge couplings explicitly
from experimental inputs at the $Z$-pole and prepare them for comparison
with the internal overlap formulation.

\subsection{Electroweak inputs}

We use the following $\overline{\mathrm{MS}}$-scheme inputs at $\mu=M_Z$:
\begin{equation}
\alpha(M_Z)^{-1} = 127.95,
\qquad
\sin^2\theta_W(M_Z) = 0.23122,
\qquad
\alpha_s(M_Z) = 0.1179.
\label{eq:EWinputs}
\end{equation}

These values are representative of current PDG averages and are sufficient
for the present analysis.

\subsection{Definition of gauge couplings}

Define the electromagnetic coupling
\begin{equation}
e := \sqrt{4\pi \alpha(M_Z)}.
\end{equation}

Then the electroweak couplings are
\begin{equation}
g_2 = \frac{e}{\sin\theta_W},
\qquad
g' = \frac{e}{\cos\theta_W},
\end{equation}
and the GUT-normalized hypercharge coupling is
\begin{equation}
g_1 = \sqrt{\frac{5}{3}}\, g'.
\end{equation}

The strong coupling is
\begin{equation}
g_3 = \sqrt{4\pi \alpha_s(M_Z)}.
\end{equation}

\subsection{Explicit numerical evaluation}

From \eqref{eq:EWinputs}:
\[
\alpha(M_Z) \approx 0.007816,
\qquad
e = \sqrt{4\pi\alpha} \approx 0.3135.
\]

Also,
\[
\sin\theta_W = \sqrt{0.23122} \approx 0.4809,
\qquad
\cos\theta_W \approx 0.8768.
\]

Thus,
\[
g_2(M_Z) = \frac{0.3135}{0.4809} \approx 0.6517,
\]
\[
g'(M_Z) = \frac{0.3135}{0.8768} \approx 0.3576,
\]
\[
g_1(M_Z) = \sqrt{\frac{5}{3}}\times 0.3576 \approx 0.4614.
\]

For QCD,
\[
g_3(M_Z) = \sqrt{4\pi\times 0.1179} \approx 1.2172.
\]

\begin{table}[h]
\centering
\begin{tabular}{lccc}
\toprule
Scale & $g_1$ (GUT) & $g_2$ & $g_3$ \\
\midrule
$M_Z$ & $0.4614$ & $0.6517$ & $1.2172$ \\
\bottomrule
\end{tabular}
\caption{Gauge couplings at $\mu=M_Z$ computed from \eqref{eq:EWinputs}.}
\end{table}

% ===========================
\section{One-loop renormalization group running to $\mu_\Theta$}
% ===========================

We now evolve the couplings to a common matching scale
\[
\mu_\Theta = 5~\mathrm{TeV},
\]
which we later identify with the minimal coherent-sector gap scale.

\subsection{RG equations}

As adopted in the MTT amplitudes construction, the one-loop RG equations are
\begin{equation}
\mu\frac{dg_a}{d\mu} = \frac{b_a}{16\pi^2} g_a^3,
\qquad
(b_1,b_2,b_3)=\left(\frac{41}{10},-\frac{19}{6},-7\right).
\label{eq:beta1loop}
\end{equation}

Equivalently,
\begin{equation}
g_a^{-2}(\mu) = g_a^{-2}(\mu_0)
- \frac{b_a}{8\pi^2}\ln\frac{\mu}{\mu_0}.
\label{eq:RGinverse}
\end{equation}

\subsection{Logarithmic factor}

Set $\mu_0=M_Z=91.1876~\mathrm{GeV}$ and $\mu=\mu_\Theta=5000~\mathrm{GeV}$.
Then
\[
L := \ln\frac{\mu_\Theta}{M_Z}
= \ln\left(\frac{5000}{91.1876}\right)
\approx \ln(54.82) \approx 3.999.
\]

Also,
\[
8\pi^2 \approx 78.9568.
\]

\subsection{Running the couplings}

Compute the shifts
\[
\Delta_a := \frac{b_a}{8\pi^2}L.
\]

Explicitly:
\[
\Delta_1 \approx \frac{4.1}{78.9568}\times 3.999 \approx 0.2077,
\]
\[
\Delta_2 \approx \frac{-3.1667}{78.9568}\times 3.999 \approx -0.1604,
\]
\[
\Delta_3 \approx \frac{-7}{78.9568}\times 3.999 \approx -0.3546.
\]

Now evolve each coupling.

\paragraph{Hypercharge.}
\[
g_1^{-2}(M_Z) = \frac{1}{0.4614^2} \approx 4.6966,
\]
\[
g_1^{-2}(\mu_\Theta) \approx 4.6966 - 0.2077 = 4.4889,
\]
\[
g_1(\mu_\Theta) \approx \frac{1}{\sqrt{4.4889}} \approx 0.4720.
\]

\paragraph{Weak SU(2).}
\[
g_2^{-2}(M_Z) = \frac{1}{0.6517^2} \approx 2.3534,
\]
\[
g_2^{-2}(\mu_\Theta) \approx 2.3534 + 0.1604 = 2.5138,
\]
\[
g_2(\mu_\Theta) \approx \frac{1}{\sqrt{2.5138}} \approx 0.6306.
\]

\paragraph{Strong SU(3).}
\[
g_3^{-2}(M_Z) = \frac{1}{1.2172^2} \approx 0.6750,
\]
\[
g_3^{-2}(\mu_\Theta) \approx 0.6750 + 0.3546 = 1.0296,
\]
\[
g_3(\mu_\Theta) \approx \frac{1}{\sqrt{1.0296}} \approx 0.9854.
\]

\begin{table}[h]
\centering
\begin{tabular}{lccc}
\toprule
Scale & $g_1$ (GUT) & $g_2$ & $g_3$ \\
\midrule
$5~\mathrm{TeV}$ & $0.4720$ & $0.6306$ & $0.9854$ \\
\bottomrule
\end{tabular}
\caption{Gauge couplings evolved to $\mu_\Theta=5~\mathrm{TeV}$ at one loop.}
\end{table}

% ===========================
\section{Numerical $\Theta$--constraints from gauge data}
% ===========================

Using the overlap relation
\[
\frac{1}{g_a^2} \propto I_a,
\]
ratios of couplings give ratios of overlap integrals:
\begin{equation}
\frac{I_b}{I_a} = \left(\frac{g_a}{g_b}\right)^2.
\end{equation}

\subsection{Explicit overlap-ratio targets}

At $\mu_\Theta=5~\mathrm{TeV}$:
\[
\frac{I_2}{I_1}
= \left(\frac{0.4720}{0.6306}\right)^2
\approx 0.560,
\]
\[
\frac{I_3}{I_1}
= \left(\frac{0.4720}{0.9854}\right)^2
\approx 0.229.
\]

\begin{equation}
\boxed{
\frac{I_2}{I_1}(\mu_\Theta)\approx 0.560,
\qquad
\frac{I_3}{I_1}(\mu_\Theta)\approx 0.229.
}
\end{equation}

These two numbers are the experimentally anchored $\Theta$–constraints
that must be satisfied by the internal geometry at the matching scale.

\begin{remark}[Target uncertainties]
The numerical targets $I_2/I_1$ and $I_3/I_1$ inherit small experimental uncertainties
from $(\alpha(M_Z),\sin^2\theta_W(M_Z),\alpha_s(M_Z))$ and small scheme systematics
from the choice of one-loop running and omission of threshold corrections.
These uncertainties can be propagated explicitly; in the present conservative core we treat
the stated targets as central values and defer a full error budget to follow-up work.
\end{remark}


\section{Transition to the geometric $\Theta$-problem}

The next step is to determine whether the internal geometry of MTT
admits overlap integrals $I_1,I_2,I_3$ satisfying the above ratios
while maintaining the baseline spectral gap $\lambda_{\ast}=0.25$.

This problem is addressed in the next section.
% ===========================
\section{Baseline internal geometry and spectral gap bounds}
% ===========================

We now specify the baseline internal geometry used in the MTT Foundation
and collect the spectral gap inequalities required for admissibility.

\subsection{Internal metric structure}

At each spacetime point $y\in Y^4$, the internal bundles have the form
\[
B_n|_y \simeq S^1_{\mathrm{cen}} \times \Sigma_n,
\qquad n=1,2,3,
\]
where:
\begin{itemize}
\item $S^1_{\mathrm{cen}}$ is a common central circle,
\item $\Sigma_1$ is an auxiliary circle,
\item $\Sigma_2$ is a lens-type surface,
\item $\Sigma_3$ is a nilmanifold leaf.
\end{itemize}

The baseline metric ansatz (Foundation, Appendix B / Section 14) is block-diagonal,
with parameters:
\begin{itemize}
\item $R_1$ — radius of $\Sigma_1$,
\item $R_{\mathrm{lens}}$ — lens radius,
\item $f_2$ — lens warp factor,
\item $h_0$ — Cheeger constant parameter controlling the nil spectrum.
\end{itemize}

\subsection{Spectral gap bounds}

The MTT Foundation establishes the following lower bounds on the smallest
nonzero eigenvalues of the internal Laplacians.

\paragraph{Circle factor $\Sigma_1$.}
For a circle of radius $R_1$,
\begin{equation}
\lambda_{\Sigma_1} \sim \frac{1}{R_1^2}.
\label{eq:gapSigma1}
\end{equation}

\paragraph{Lens factor $\Sigma_2$.}
For the lens sheet with effective radius $f_2R_{\mathrm{lens}}$,
\begin{equation}
\lambda_{\mathrm{lens}} \ge \frac{2}{(f_2R_{\mathrm{lens}})^2}.
\label{eq:gapLens}
\end{equation}

\paragraph{Nil factor $\Sigma_3$.}
For the nil leaf, the spectral gap is controlled by the Cheeger constant:
\begin{equation}
h \ge h_0 > 0
\quad\Rightarrow\quad
\lambda_{\mathrm{nil}} \ge \frac{h_0^2}{4}.
\label{eq:gapNil}
\end{equation}

\subsection{Baseline numerical values}

In the worked numerical baseline of the Foundation,
\begin{equation}
h_0 = 1
\quad\Rightarrow\quad
\lambda_{\mathrm{nil}} = \frac{1}{4} = 0.25.
\end{equation}

We enforce the \emph{baseline lower bound} $\lambda_{\ast}\ge 0.25$ using $h_0=1$.
This sets a conservative admissibility floor; explicit realizations of the nil sector may
exceed this bound and need not saturate it.


% ===========================
\section{Geometric overlap model}
% ===========================

We now introduce a minimal, computable overlap model consistent with
the amplitudes construction.

\subsection{Factorized overlap ansatz}

Assume the harmonic representatives factorize as
\[
\omega_a(\theta,u) = \omega_{S^1}(\theta)\,\phi_a(u),
\qquad u\in\Sigma_a,
\]
with
\[
\|\omega_{S^1}\|_{L^2(S^1_{\mathrm{cen}})} = 1,
\qquad
\|\phi_a\|_{L^2(\Sigma_a)} = 1.
\]

Then the overlap integrals reduce to
\begin{equation}
I_a = \int_{S^1_{\mathrm{cen}}\times\Sigma_a}
|\omega_{S^1}|^2\,|\phi_a|^2\,d\mu
= \mathrm{Vol}(\Sigma_a).
\label{eq:IasVolume}
\end{equation}

\subsection{Explicit volume expressions}

We parameterize the volumes as:
\begin{align}
I_1 &= 2\pi R_1,
\label{eq:I1}\\
I_2 &= \kappa_\ell (f_2R_{\mathrm{lens}})^2,
\label{eq:I2}\\
I_3 &= \kappa_n\,\mathcal{S}_n,
\label{eq:I3}
\end{align}
where:
\begin{itemize}
\item $\kappa_\ell$ is the unit-area normalization of the lens sheet,
\item $\kappa_n\,\mathcal{S}_n$ is the effective nil overlap weight,
\item $\mathcal{S}_n$ is not assumed equal to a simple metric area.
\end{itemize}

This separation is required because the nil spectral gap
\eqref{eq:gapNil} is controlled by $h_0$, not directly by metric scale.

% ===========================
\section{Solving the $\Theta$--system}
% ===========================

We now solve the overlap constraints from Section~4 of Part~II
together with the spectral gap bounds.

\subsection{Overlap constraints}

From \eqref{eq:targets} and \eqref{eq:I1}--\eqref{eq:I3}:
\begin{align}
\frac{I_2}{I_1}
&=
\frac{\kappa_\ell (f_2R_{\mathrm{lens}})^2}{2\pi R_1}
= 0.560,
\label{eq:Theta1}\\
\frac{I_3}{I_1}
&=
\frac{\kappa_n\mathcal{S}_n}{2\pi R_1}
= 0.229.
\label{eq:Theta2}
\end{align}

Solving \eqref{eq:Theta1} gives
\begin{equation}
(f_2R_{\mathrm{lens}})^2
=
\frac{0.560\cdot 2\pi}{\kappa_\ell}\,R_1
=
\frac{3.518}{\kappa_\ell}\,R_1,
\label{eq:RlSolve}
\end{equation}
hence
\begin{equation}
R_{\mathrm{lens}}
=
\frac{\sqrt{(3.518/\kappa_\ell)\,R_1}}{f_2}.
\label{eq:RlFinal}
\end{equation}

From \eqref{eq:Theta2}:
\begin{equation}
\kappa_n\mathcal{S}_n
=
0.229\cdot 2\pi\,R_1
=
1.439\,R_1.
\label{eq:SnSolve}
\end{equation}

\subsection{Imposing gap inequalities}

We require:
\begin{align}
\lambda_{\Sigma_1} &\ge 0.25,
\label{eq:gap1}\\
\lambda_{\mathrm{lens}} &\ge 0.25,
\label{eq:gap2}\\
\lambda_{\mathrm{nil}} &\ge 0.25.
\label{eq:gap3}
\end{align}

\paragraph{Nil.}
Set $h_0=1$, so \eqref{eq:gapNil} gives $\lambda_{\mathrm{nil}}=0.25$.

\paragraph{Circle.}
From \eqref{eq:gapSigma1} and \eqref{eq:gap1},
\[
\frac{1}{R_1^2}\ge 0.25
\quad\Rightarrow\quad
R_1\le 2.
\]

\paragraph{Lens.}
From \eqref{eq:gapLens} and \eqref{eq:gap2},
\[
\frac{2}{(f_2R_{\mathrm{lens}})^2}\ge 0.25
\quad\Rightarrow\quad
(f_2R_{\mathrm{lens}})^2\le 8.
\]

Substitute \eqref{eq:RlSolve}:
\[
\frac{3.518}{\kappa_\ell}\,R_1 \le 8
\quad\Rightarrow\quad
R_1 \le \frac{8\kappa_\ell}{3.518}\approx 2.274\,\kappa_\ell.
\]

\subsection{Existence theorem}

\begin{theorem}[Existence of an admissible $\Theta$--region]
Fix $\mu_\Theta=5~\mathrm{TeV}$ and choose
\[
h_0=1,\qquad \kappa_\ell=1,\qquad f_2\ge 1.
\]
Then for every $R_1\in(0,2]$ there exists a choice of
$R_{\mathrm{lens}}$ and $\mathcal{S}_n$ given by
\eqref{eq:RlFinal} and \eqref{eq:SnSolve} such that:
\begin{enumerate}
\item the gauge overlap constraints
$\frac{I_2}{I_1}\approx 0.560$ and $\frac{I_3}{I_1}\approx 0.229$
are satisfied;
\item the spectral gap inequalities
$\lambda_{\Sigma_1},\lambda_{\mathrm{lens}},\lambda_{\mathrm{nil}}\ge 0.25$
hold;
\item the minimal spectral gap remains $\lambda_{\ast}=0.25$.
\end{enumerate}
\end{theorem}

\begin{proof}
Let $R_1\in(0,2]$. Then $\lambda_{\Sigma_1}\ge 0.25$.
With $\kappa_\ell=1$ and $f_2\ge 1$, \eqref{eq:RlFinal} implies
$(f_2R_{\mathrm{lens}})^2=3.518R_1\le 7.036<8$, hence
$\lambda_{\mathrm{lens}}>0.25$.
With $h_0=1$, $\lambda_{\mathrm{nil}}=0.25$.
The overlap constraints are satisfied by construction.
\end{proof}

\section{Interpretation}

The existence theorem shows that Standard Model gauge couplings
do not conflict with the MTT coherence and gap requirements.
The nil sector saturates the minimal gap and therefore controls
the onset of noncoherent physics.

\begin{remark}
The quantities $\kappa_n$ and $\mathcal{S}_n$ encode nontrivial
nilmanifold structure and are not free once explicit harmonic
representatives are specified.
\end{remark}

\begin{remark}[Coherent sector stability]
All results in this work are conditional on the stability of the coherent
sector.
If the coherent sector were dynamically unstable, the theory would not merely
predict different numerical values; rather, it would fail to support persistent
observables, reproducible measurements, or probabilistic interpretation.
Accordingly, coherent sector stability is treated here as a prerequisite for
physics rather than a phenomenological assumption.
\end{remark}

% ===========================
\section{Massless gauge-sector normalization via the twistor corner}
% ===========================

The remaining technical ambiguity in the $\Theta$--closure analysis concerns
the explicit definition and normalization of the nonabelian overlap integrals
$I_2$ and $I_3$ associated with the $SU(2)$ and $SU(3)$ gauge sectors.
In this section we show that these quantities admit a canonical and
gap--controlled definition in the high--coherence (twistor) corner of
Modal Triplet Theory.

\begin{remark}[Twistor corner as a computational device]
In this paper the twistor corner is used solely as a \emph{computational normalization regime}
for the massless coherent gauge sector, providing canonical harmonic representatives and
gap-controlled error bounds. No claim is made that Nature resides exactly in this corner,
nor that twistor space is fundamental.
\end{remark}


\subsection{High--coherence regime and admissible truncation}

The twistor formulation of MTT establishes the existence of a high--coherence
regime in which:
\begin{enumerate}
\item the spectral gap $\lambda_Q$ separating coherent and noncoherent modes
is parametrically large;
\item the effective generator admits a Schur--Feshbach reduction onto a
finite--dimensional coherent subspace;
\item truncation errors are bounded by explicit operator--norm estimates
of order $O(\lambda_Q^{-1})$.
\end{enumerate}

\begin{assumption}[Twistor-corner admissibility]
We assume that the $\Theta$--matching scale $\mu_\Theta$ lies within the
high--coherence regime for the massless gauge sector, so that the
Schur--Feshbach reduction is valid with remainder bounded by
$O(\lambda_Q^{-1})$.
\end{assumption}

\begin{remark}
This assumption invokes no ultraviolet completion and no string-theoretic
structure; it relies only on the coherence and gap hypotheses already stated
in Sections~2 and~10.
\end{remark}

\subsection{Canonical definition of the massless gauge subspace}

Let $\mathcal H$ denote the full internal Hilbert space of modal fields.
Define $\mathcal H_0\subset\mathcal H$ as the coherent massless gauge subspace
selected by the twistor corner, characterized by:
\begin{enumerate}
\item self--duality of the Yang--Mills curvature;
\item vanishing effective mass under the reduced generator;
\item holomorphic encoding via the Ward correspondence.
\end{enumerate}

\begin{definition}[Twistor-corner massless sector]
$\mathcal H_0$ is defined as the image of the coherent projector $P_0$
associated with the self--dual Yang--Mills corner of the MTT configuration
space.
\end{definition}

This definition is canonical within the admissible slab and does not depend
on additional choices.

\subsection{Schur--Feshbach reduction and error control}

Let $L$ be the full internal generator and write
\[
L =
\begin{pmatrix}
P_0 L P_0 & P_0 L Q \\
Q L P_0 & Q L Q
\end{pmatrix},
\qquad Q=1-P_0.
\]

The effective generator on $\mathcal H_0$ is given by the Schur--Feshbach map
\[
L_{\mathrm{eff}}
=
P_0 L P_0
-
P_0 L Q (Q L Q)^{-1} Q L P_0.
\]

The twistor analysis proves the operator bound
\[
\|P_0 L Q (Q L Q)^{-1} Q L P_0\|
\;\le\;
C\,\lambda_Q^{-1},
\]
for some constant $C$ independent of infrared data.

\subsection{Harmonic representatives and normalization}

\begin{definition}[Twistor-corner harmonic representatives]
For $a=2,3$, let $\omega_a^{(0)}$ denote the unique $L^2$--normalized
harmonic representative in $\mathcal H_0$ associated with the $SU(a)$
gauge sector, defined as the minimizer of the internal $L^2$ norm within
the corresponding twistor cohomology class and satisfying
\[
\|\omega_a^{(0)}\|_{L^2(B_a)} = 1.
\]
\end{definition}

This variational characterization fixes normalization uniquely, up to
gauge transformations.

\begin{remark}
The abelian overlap $I_1$ remains fixed by the explicit $S^1_{\mathrm{cen}}$
normalization of Appendix~G and is not modified by the twistor construction.
\end{remark}

\subsection{Leading--order overlap integrals}

\begin{definition}[Leading--order overlaps]
The leading--order nonabelian overlap integrals are defined by
\[
I_a^{(0)} := \int_{B_a}
\langle \omega_a^{(0)}, \omega_a^{(0)} \rangle \, d\mu_{B_a},
\qquad a=2,3.
\]
\end{definition}

These depend only on the massless coherent geometry encoded by the twistor
corner.

\subsection{Control of overlap deviations}

Let $I_a$ denote the full overlap integrals defined in Section~4.

\begin{lemma}[Overlap stability under coherent truncation]
There exist constants $C_a>0$ such that
\[
|I_a - I_a^{(0)}| \;\le\; C_a\,\lambda_Q^{-1},
\qquad a=2,3.
\]
\end{lemma}

\begin{proof}
By Kato--Rellich perturbation theory, spectral projections and their
associated $L^2$ norms depend continuously on the generator.
The Schur--Feshbach bound controls the perturbation of $L$ on $\mathcal H_0$,
implying the stated estimate.
\end{proof}

\subsection{Twistor-corner $\Theta$--closure at leading order}

Combining the above with the experimentally extracted targets at
$\mu_\Theta=5~\mathrm{TeV}$,
\[
\frac{I_2}{I_1} \approx 0.560,
\qquad
\frac{I_3}{I_1} \approx 0.229,
\]
we obtain:

\begin{theorem}[Twistor-corner $\Theta$--closure condition]
If the twistor-corner admissibility assumption holds at $\mu_\Theta$, then
$\Theta$--closure at leading order reduces to the computable conditions
\[
\frac{I_2^{(0)}}{I_1} \approx 0.560,
\qquad
\frac{I_3^{(0)}}{I_1} \approx 0.229,
\]
with controlled corrections of order $O(\lambda_Q^{-1})$.
\end{theorem}

\begin{remark}
Failure of these relations falsifies either the baseline internal geometry
or the twistor-corner admissibility hypothesis.
\end{remark}

\subsection{Scope}

This section does not claim that the twistor corner describes the full
physical gauge sector. It provides a canonical, gap--controlled
computational corner in which $I_2$ and $I_3$ are well defined and
numerically accessible.

% ===========================
\section{Absolute scale calibration and quantum--gravity consistency}
% ===========================

Parts~I--III established a nonempty admissible $\Theta$--region in
dimensionless internal parameters.
We now connect this region to physical units by a single calibration choice,
and derive consequences for quantum gravity and early--universe physics.

\subsection{Calibration of the minimal gap scale}

From the MTT Foundation, the minimal coherent--sector gap is
\[
\lambda_{\ast} = 0.25.
\]

To assign physical units, we introduce a calibration assumption.

\begin{assumption}[Gap--scale calibration]
We identify the minimal coherent--sector excitation scale with the matching scale
\[
E_{\mathrm{gap,min}} := \mu_\Theta = 5~\mathrm{TeV}.
\]
\end{assumption}

This choice is conservative: it lies above current collider bounds
and below any Planckian or GUT-scale assumptions.

\begin{remark}
The MTT formalism itself does not fix this identification; it only requires
$\lambda_{\ast}>0$. The calibration is therefore explicit and falsifiable.
\end{remark}

\subsection{Physical length scale}

Eigenvalues of the internal Laplacian scale as $\lambda \sim R^{-2}$.
Thus, identifying $E_{\mathrm{gap,min}}=\sqrt{\lambda_{\ast}}/R_0$ gives
\begin{equation}
\frac{1}{R_0}
=
\frac{E_{\mathrm{gap,min}}}{\sqrt{\lambda_{\ast}}}
=
\frac{5~\mathrm{TeV}}{0.5}
=
10~\mathrm{TeV}.
\end{equation}

Hence
\begin{equation}
R_0 = (10~\mathrm{TeV})^{-1}.
\end{equation}

Using $1~\mathrm{GeV}^{-1}\approx 1.97\times 10^{-16}\,\mathrm{m}$,
\begin{equation}
R_0 \approx 10^{-4}~\mathrm{GeV}^{-1}
\approx 2\times 10^{-20}\,\mathrm{m}.
\end{equation}

This is the characteristic physical size associated with the minimal
coherent--sector gap.

% ===========================
\section{SPT proper--time gap and UV damping}
% ===========================

The MTT quantum--gravity construction introduces a spectral proper--time (SPT)
filter with a positive lower bound $\tau_0>0$ on its support.

\subsection{SPT damping scale}

The filter produces Gaussian damping of graviton propagators of the form
\[
e^{-\tau_0 k^2}.
\]

Dimensional consistency suggests identifying the damping scale as
\[
\Lambda := \tau_0^{-1/2}.
\]

\begin{assumption}[Minimal SPT identification]
We identify
\[
\Lambda \equiv \tau_0^{-1/2} \sim E_{\mathrm{gap,min}} = 5~\mathrm{TeV}.
\]
\end{assumption}

\begin{remark}
The constructive QG results require only $\tau_0>0$.
The numerical identification of $\tau_0$ follows from the chosen gap calibration.
\end{remark}

Thus
\[
\tau_0 \sim (5~\mathrm{TeV})^{-2}.
\]

% ===========================
\section{Infrared consistency of gravity}
% ===========================

We verify that the calibrated scale $\Lambda\sim 5~\mathrm{TeV}$ does not
conflict with known infrared tests of gravity.

\subsection{Size of higher--derivative corrections}

Assume the leading gravity--sector corrections appear as
\[
\mathcal{O}\!\left(\frac{\Box}{\Lambda^2}\right).
\]

For a process with characteristic frequency $\omega$,
the relative size scales as
\[
\left(\frac{\omega}{\Lambda}\right)^2.
\]

\paragraph{Gravitational waves today.}
LIGO--band frequencies satisfy $\omega\sim 10^2$--$10^3~\mathrm{Hz}$.
Using
\[
1~\mathrm{Hz}\approx 4.14\times 10^{-15}~\mathrm{eV},
\qquad
\Lambda=5\times 10^{12}~\mathrm{eV},
\]
we estimate
\begin{equation}
\left(\frac{\omega}{\Lambda}\right)^2
\lesssim
\left(\frac{10^3\times 4\times 10^{-15}}{5\times 10^{12}}\right)^2
\sim 10^{-48}.
\end{equation}

Thus any modification of GW propagation is entirely negligible.

\subsection{Higher--curvature sanity check}

A representative EFT correction is an $R^2$ term,
equivalent to a scalar mode with mass $m_0$.
Identifying $m_0\sim\Lambda$ gives a Yukawa range
\begin{equation}
\lambda_0 = m_0^{-1} \sim (5~\mathrm{TeV})^{-1}
\approx 4\times 10^{-20}\,\mathrm{m}.
\end{equation}

This range lies many orders of magnitude below the reach of current
inverse--square--law or fifth--force experiments.

\begin{remark}
Large dimensionless coefficients of $R^2$ are expected because
$m_0\ll M_{\mathrm{Pl}}$; this does not signal strong gravity in the IR.
\end{remark}

% ===========================
\section{Inflationary regimes and inadmissibility}
% ===========================

We now examine early--universe regimes in light of the calibrated
coherence scale $\Lambda$.

\subsection{Inflationary tensor scale}

In conventional slow--roll inflation, the tensor amplitude is controlled by
the Hubble scale $H$,
\[
A_t \sim \frac{H^2}{M_{\mathrm{Pl}}^2}.
\]

Current bounds allow $H\sim 10^{13}$--$10^{14}~\mathrm{GeV}$.

\subsection{Comparison with the coherence scale}

With $\Lambda\sim 5\times 10^3~\mathrm{GeV}$,
\begin{equation}
\left(\frac{H}{\Lambda}\right)^2
\sim
\left(\frac{10^{13}\text{--}10^{14}}{5\times 10^3}\right)^2
\sim 10^{18}\text{--}10^{20}.
\end{equation}

Thus $H\gg\Lambda$ in standard high--scale inflation.

\begin{theorem}[Inflationary inadmissibility]
Any inflationary phase with $H\gg\Lambda$ is inadmissible in Modal Triplet Theory:
the coherent--sector truncation fails, and quantum perturbation theory is not
well defined.
\end{theorem}

\begin{proof}
Admissibility requires that higher--derivative corrections remain parametrically
small. For $H\gg\Lambda$, terms of order $H^2/\Lambda^2$ dominate, violating
the spectral gap and bounded--projection assumptions.
\end{proof}

\begin{corollary}
Standard high--scale inflationary tensor predictions are not physical
observables under the calibration $\Lambda\sim 5~\mathrm{TeV}$.
\end{corollary}

\begin{remark}
Low--scale inflation with $H\ll\Lambda$ remains admissible but predicts
negligible primordial tensor amplitudes.
\end{remark}

% ===========================
\section{Falsifiability}
% ===========================

The $\Theta$--closure program yields explicit experimental discriminators.

\subsection{Gauge sector}

Failure to realize overlap ratios
\[
\frac{I_2}{I_1}\approx 0.560,
\qquad
\frac{I_3}{I_1}\approx 0.229
\]
within the baseline internal geometry falsifies the assumed
bundle structure or normalization.

\subsection{Cosmology}

A robust detection of primordial tensor modes requiring $H\gg\Lambda$
would falsify either:
\begin{enumerate}
\item the gap--scale calibration, or
\item the coherence assumptions of MTT.
\end{enumerate}

\subsection{Quantum gravity}

Observation of unsuppressed high--energy graviton scattering
or violation of Gaussian UV damping would contradict the SPT gap assumption.

\section{Conclusion}

In this paper we have established a conservative and fully explicit
$\Theta$--closure framework in Modal Triplet Theory.

Beginning with experimentally measured Standard Model gauge couplings, we
derived numerical $\Theta$--targets for the nonabelian overlap ratios
$I_2/I_1$ and $I_3/I_1$ at a common matching scale
$\mu_\Theta=5~\mathrm{TeV}$.
We formulated admissibility conditions on the internal geometry in terms of
spectral gap and bounded projection requirements and proved the existence of a
nonempty $\Theta$--region satisfying these constraints.

A key technical issue—the normalization of nonabelian harmonic representatives—
was isolated and addressed by introducing a high--coherence (twistor) corner,
which provides canonical massless gauge modes with controlled corrections
suppressed by the coherence gap.
This construction does not assume that Nature resides exactly in the twistor
corner; rather, it supplies a well--defined computational regime within which
$\Theta$--closure can be formulated precisely.

The present paper does not claim to compute the overlap integrals explicitly.
Instead, it defines the $\Theta$--targets, the admissibility conditions, and the
normalization framework required for such computations.
In Paper~II, these targets are realized explicitly by direct geometric
evaluation of the overlaps on the baseline internal space
$S^1_{\mathrm{cen}}\times L(3,1)\times\Gamma\backslash\mathrm{Nil}_3$
(Route~A).
In Paper~III, the same overlaps are derived independently using twistor--action
normalization (Route~B), providing a nontrivial cross--check.

Taken together, these results show that the leading--order nonabelian gauge
overlaps in MTT are not free parameters but are fixed—up to controlled
corrections—by internal structure and coherence conditions.
This completes the conservative $\Theta$--closure of the gauge sector at leading
order and establishes a solid foundation for extending the program to
gravitational and cosmological observables.

% ===========================
\section*{References}
% ===========================

\begin{enumerate}
\item Particle Data Group (PDG), Review of Particle Physics (for $M_Z$, $\alpha(M_Z)$, $\sin^2\theta_W(M_Z)$, $\alpha_s(M_Z)$).
\item M. E. Machacek and M. T. Vaughn, ``Two-loop renormalization group equations in a general quantum field theory,'' Nucl. Phys. B222 (1983) 83--103.  % for RGE conventions
\item T. Kato, \emph{Perturbation Theory for Linear Operators}, Springer (1995). % for spectral projection continuity
\item Peter Nero, \emph{Modal Triplet Theory: Foundation} (MTT corpus).
\item Peter Nero, \emph{Modal Triplet Theory: Quantum Amplitudes from Modal Geometry} (MTT corpus).
\item Peter Nero, \emph{Twistor Encodings as High-Coherence Limits of Modal Triplet Theory} (MTT corpus).
\end{enumerate}

\end{document}

